\chapter*{مقدمة}
	التحليل العقدي هو فرع من فروع الرياضيات الذي يختص بدراسة الدوال التي تُعرّف على الأعداد العقدية. حيث أن الأعداد العقدية هي أعداد تتكون من جزئين: جزء حقيقي وآخر تخيلي. 
	
	يعد هذا المجال أساسيًا لفهم سلوك الدوال في المجال العقدي وكيفية تأثير الأجزاء الحقيقية والتخيلية للأعداد العقدية على تلك الدوال. ومن خلال دراسة هذه الدوال، يمكننا فهم خصائص معقدة تتعلق بالتحليل الرياضي، بالإضافة إلى تطبيقات عملية في العديد من المجالات مثل الفيزياء والهندسة.
	
	تعد التطبيقات العملية للتحليل العقدي واسعة جدًا، فهي تستخدم في العديد من الظواهر الطبيعية مثل تفاعلات الجسيمات في فيزياء الكم، وأيضًا في تحليل الدوائر الكهربائية والنماذج الرياضية في الهندسة. يعتبر التحليل العقدي أداة قوية لفهم العديد من الأنظمة الرياضية المعقدة، كما أنه يعد ركيزة أساسية لحل المعادلات التفاضلية التي تظهر في التطبيقات الفيزيائية والهندسية.
	
	من خلال هذا المجال، يتمكن الباحثون من دراسة التكامل العقدي ودوال أخرى في المستوى العقدي، مما يعزز القدرة على التعامل مع النماذج الرياضية المعقدة وفهم الديناميكيات الدقيقة التي تحكم الأنظمة المختلفة.
	
